% 2010 年江西卷文科第 12 题：四个图象选项（A）—（D）。
% 来源 PDF：普通高考/2010/2010江西文.pdf，SHA-256 c8c86a7272df02211069e27098dcaf04a5cd8af9a9b4d54ffd14a8eb91347af3。
% 原页：第 1 页（A3 三栏）第 3 栏中部，书页码 609。
% 题设三条曲线：y=sin2x（周期 π）、y=sin(x+π/6)、y=sin(x−π/3)（后两条周期 2π，
%   相位差 π/2）；前者的频率是后两条的两倍，同区间内波峰数约为两倍。
% 重绘依据：原页每幅仅约 4cm 宽、曲线在上边框被裁切，逐像素判读不可靠，故本图为
%   按题设关系绘制的示意图：A、B、D 三幅中 sin2x 呈二倍频、两条慢曲线相差四分之一
%   周期；C 幅三条曲线同频（无二倍频），即错误项。已在 PROGRESS 登记。
% 实现说明：TikZ 的 plot 以度为角度单位，\x 取 0..720 度并线性映射到面板宽度，
%   正弦用 (\x) 度计算，避免把弧度误当角度。
\begin{tikzpicture}[semithick]
  \def\w{4.6}       % 面板宽 (cm)
  \def\hh{1.5}      % 面板高 (cm)
  \def\amp{0.62}    % 振幅 (cm)
  \def\xscale{0.0063889}  % cm/度 = 4.6/720
  % 正确画法：#1 为整体相位平移（度）
  \newcommand{\goodpanel}[1]{%
    % y = sin 2x ：720 度内四个波峰
    \draw[thick] plot[domain=0:720,samples=361]
      ({\x*\xscale},{\amp*sin(2*\x+#1)});
    % y = sin(x + pi/6)：三十度为 30 度相位
    \draw[thick] plot[domain=0:720,samples=241]
      ({\x*\xscale},{\amp*sin(\x+30+#1)});
    % y = sin(x - pi/3)：-60 度相位，与上一条相差 90 度
    \draw[thick] plot[domain=0:720,samples=241]
      ({\x*\xscale},{\amp*sin(\x-60+#1)});
  }
  % 错误画法：三条曲线同为慢频（无二倍频），相位各不相同
  \newcommand{\badpanel}{%
    \draw[thick] plot[domain=0:720,samples=241]
      ({\x*\xscale},{\amp*sin(\x+20)});
    \draw[thick] plot[domain=0:720,samples=241]
      ({\x*\xscale},{\amp*sin(\x+110)});
    \draw[thick] plot[domain=0:720,samples=241]
      ({\x*\xscale},{\amp*sin(\x+200)});
  }
  % (A)
  \begin{scope}
    \draw (0,-\hh/2) rectangle (\w,\hh/2);
    \goodpanel{90}
    \node at (\w/2,-\hh/2-0.35) {（A）};
  \end{scope}
  % (B)
  \begin{scope}[shift={(5.5,0)}]
    \draw (0,-\hh/2) rectangle (\w,\hh/2);
    \goodpanel{150}
    \node at (\w/2,-\hh/2-0.35) {（B）};
  \end{scope}
  % (C) 错误项
  \begin{scope}[shift={(0,-2.5)}]
    \draw (0,-\hh/2) rectangle (\w,\hh/2);
    \badpanel
    \node at (\w/2,-\hh/2-0.35) {（C）};
  \end{scope}
  % (D)
  \begin{scope}[shift={(5.5,-2.5)}]
    \draw (0,-\hh/2) rectangle (\w,\hh/2);
    \goodpanel{210}
    \node at (\w/2,-\hh/2-0.35) {（D）};
  \end{scope}
\end{tikzpicture}
